% =============================================================================
\chapter{API Reference}
\label{ch:api}
% =============================================================================

This chapter documents every public class, method, and function in the
toolkit.  Private and internal methods are omitted; consult the source code
for implementation details.

% -----------------------------------------------------------------------------
\section{\texttt{StructureGeometry}}
\label{sec:api:geometry}
% -----------------------------------------------------------------------------

Stores topology, material properties, boundary conditions, and DOF mappings.
Properties that depend on thickness or node positions are kept in sync through
dedicated setter methods.

\subsection{Constructor}

\begin{lstlisting}
geom = StructureGeometry(nodes, connectivity, props, constraints)
\end{lstlisting}

\begin{table}[H]
\centering
\caption{\texttt{StructureGeometry} constructor inputs.}
\begin{tabularx}{\textwidth}{llX}
\toprule
\textbf{Argument} & \textbf{Size} & \textbf{Description} \\
\midrule
\texttt{nodes}        & $N\times2$  & Node coordinates $[x,y]$ in metres. \\
\texttt{connectivity} & $M\times2$  & Element connectivity $[\text{start},\,\text{end}]$. \\
\texttt{props}        & struct      & Material/section properties (see below). \\
\texttt{constraints}  & $C\times3$  & Boundary conditions $[\text{node},\,\text{dof},\,\text{value}]$. \\
\bottomrule
\end{tabularx}
\end{table}

The \texttt{props} struct accepts the following fields.  Scalar values are
automatically expanded to $[M\times1]$ vectors.

\begin{table}[H]
\centering
\caption{\texttt{props} struct fields.}
\begin{tabularx}{\textwidth}{llcX}
\toprule
\textbf{Field} & \textbf{Units} & \textbf{Required} & \textbf{Description} \\
\midrule
\texttt{E}         & Pa       & Yes & Elastic modulus. \\
\texttt{thickness} & m        & Yes & Element thickness. \\
\texttt{width}     & m        & Yes & Out-of-plane width. \\
\texttt{ultimate}  & Pa       & No  & Ultimate tensile strength (default 0). \\
\texttt{rho}       & kg/m\textsuperscript{3} & No  & Density (default 7850). \\
\texttt{is\_rigid} & logical  & No  & Rigid element flag (default \texttt{false}). \\
\bottomrule
\end{tabularx}
\end{table}

\subsection{Mutators}

\begin{lstlisting}
geom.updateThickness(new_t)       % [Mx1] thickness vector (m)
geom.updateNodePositions(new_xy)  % [Nx2] node coordinates (m)
\end{lstlisting}

Both methods automatically recompute all dependent quantities ($A$, $I$, $L$,
$\theta_0$).

\subsection{Key Query Methods}

\begin{lstlisting}
[n1, n2]   = geom.getMemberNodes(id)   % node IDs for element id
dofs       = geom.getMemberDOF(id)     % [6x1] global DOF indices
xy         = geom.getNodeCoords(id)    % [1x2] coordinates
dofs       = geom.getNodeDOF(id)       % [3x1] global DOF indices
tf         = geom.isDOFConstrained(d)  % true if DOF is fixed
mass       = geom.getTotalMass()       % total structure mass (kg)
info       = geom.getStructureInfo()   % summary struct
\end{lstlisting}

% -----------------------------------------------------------------------------
\section{\texttt{ElementMatrices}}
\label{sec:api:matrices}
% -----------------------------------------------------------------------------

Manages pre-computed stiffness and transformation matrices for every element.
Uses dirty flags for lazy recomputation.

\subsection{Constructor}

\begin{lstlisting}
em = ElementMatrices(geometry)
\end{lstlisting}

Computes all local, transformation, and global matrices from the initial
geometry.

\subsection{Key Methods}

\begin{lstlisting}
% Retrieval (lazy recomputation)
k   = em.getLocalMatrix(id)        % [6x6] local stiffness
K   = em.getGlobalMatrix(id)       % [6x6] global stiffness
T   = em.getTransformation(id)     % [6x6] transformation matrix
th  = em.getCurrentAngle(id)       % current orientation (rad)
L   = em.getCurrentLength(id)      % current length (m)

% Geometry updates (Newton-Raphson)
em.updateMemberGeometry(id, new_L, new_theta)
em.updateAllMemberGeometry(lengths, thetas)

% Assembly
K_global = em.assembleGlobalStiffnessMatrix(geometry)

% Reset
em.resetToInitialConfiguration()

% Rigid stiffness
em.setRigidStiffnessFactor(1e4)    % change penalty multiplier
\end{lstlisting}

% -----------------------------------------------------------------------------
\section{\texttt{AnalysisState}}
\label{sec:api:state}
% -----------------------------------------------------------------------------

Mutable container for displacement, force, and residual vectors during
analysis.

\subsection{Constructor}

\begin{lstlisting}
state = AnalysisState(geometry)     % zero-initialised
\end{lstlisting}

\subsection{Key Properties}

\begin{table}[H]
\centering
\caption{\texttt{AnalysisState} key properties.}
\begin{tabularx}{\textwidth}{llX}
\toprule
\textbf{Property} & \textbf{Size} & \textbf{Description} \\
\midrule
\texttt{u}             & $D\times1$ & Full displacement vector. \\
\texttt{f\_external}   & $D\times1$ & External force vector. \\
\texttt{f\_internal}   & $D\times1$ & Internal (restoring) force vector. \\
\texttt{u\_free}       & $F\times1$ & Free-DOF displacements. \\
\texttt{residual}      & $F\times1$ & $\vect{f}_{\text{ext}} - \vect{f}_{\text{int}}$ on free DOFs. \\
\texttt{residual\_norm}& scalar     & $\|\text{residual}\|$. \\
\texttt{member\_forces\_local} & $M\times6$ & Element end forces (local). \\
\bottomrule
\end{tabularx}
\end{table}

\subsection{Key Methods}

\begin{lstlisting}
state.reset()                          % zero everything
state.copyFrom(other_state)            % deep copy
state.updateDisplacements(u_free)      % set free DOFs, map to full
state.computeInternalForces(em)        % assemble f_int from K*u
state.computeResidual()                % R = f_ext - f_int
state.applyIncrement(delta_u)          % u += delta_u
state.extractMemberDisplacements(em)   % populate member_displ_*
state.computeMemberForces(em)          % populate member_forces_*
d = state.getMaxDisplacement()         % max translational displacement
\end{lstlisting}

\begin{warnbox}[Reaction forces]
The \texttt{f\_internal} vector has fixed-DOF entries zeroed after assembly
to support the Newton--Raphson residual computation.  To obtain reaction
forces, compute $\vect{R} = \mat{K}\vect{u} - \vect{f}_{\text{ext}}$ on
the fixed DOFs after convergence.
\end{warnbox}

% -----------------------------------------------------------------------------
\section{\texttt{FEAAnalysis}}
\label{sec:api:analysis}
% -----------------------------------------------------------------------------

Main analysis controller.

\subsection{Constructor}

\begin{lstlisting}
analysis = FEAAnalysis(geometry)
analysis = FEAAnalysis(geometry, options)   % override settings
\end{lstlisting}

\subsection{Solver Settings}

\begin{table}[H]
\centering
\caption{\texttt{FEAAnalysis} solver settings (public, writable).}
\begin{tabularx}{\textwidth}{llcX}
\toprule
\textbf{Property} & \textbf{Type} & \textbf{Default} & \textbf{Description} \\
\midrule
\texttt{verbose}               & logical & \texttt{true} & Print iteration info. \\
\texttt{max\_iterations}       & int     & 1000          & Max N-R iterations per step. \\
\texttt{force\_tolerance}      & double  & $10^{-2}$     & Normalised force residual. \\
\texttt{displacement\_tolerance} & double & $10^{-6}$    & Normalised displacement increment. \\
\texttt{energy\_tolerance}     & double  & $10^{-6}$     & $|\vect{R}\T\Delta\vect{u}|$. \\
\bottomrule
\end{tabularx}
\end{table}

\subsection{Analysis Methods}

\begin{lstlisting}
results = analysis.runLinearAnalysis(load_case)
results = analysis.solveLargeDeflection(load_case, num_steps)
stress  = analysis.performStressAnalysis()
\end{lstlisting}

\subsection{Results Structs}

The \texttt{results} struct returned by both analysis methods contains:

\begin{table}[H]
\centering
\caption{Fields in the analysis results struct.}
\begin{tabularx}{\textwidth}{lX}
\toprule
\textbf{Field} & \textbf{Description} \\
\midrule
\texttt{displacements}      & $[D\times1]$ full displacement vector. \\
\texttt{displacements\_free} & $[F\times1]$ free-DOF displacements. \\
\texttt{max\_displacement}  & Maximum translational displacement magnitude. \\
\texttt{external\_forces}   & $[D\times1]$ applied force vector. \\
\texttt{nodes\_deformed}    & $[N\times2]$ deformed node coordinates. \\
\bottomrule
\end{tabularx}
\end{table}

The \texttt{stress} struct from \texttt{performStressAnalysis} contains:

\begin{table}[H]
\centering
\caption{Fields in the stress results struct.}
\begin{tabularx}{\textwidth}{lX}
\toprule
\textbf{Field} & \textbf{Description} \\
\midrule
\texttt{stresses.axial}       & $[M\times2]$ axial stress at each node (Pa). \\
\texttt{stresses.bending}     & $[M\times2]$ bending stress (Pa). \\
\texttt{stresses.shear}       & $[M\times2]$ shear stress (Pa). \\
\texttt{stresses.von\_mises}  & $[M\times2]$ von Mises stress (Pa). \\
\texttt{stresses.max\_stress} & $[M\times1]$ max von Mises per element (Pa). \\
\texttt{stresses.critical\_member} & Element ID with highest stress. \\
\texttt{FOS.member\_FOS}      & $[M\times1]$ factor of safety per element. \\
\texttt{FOS.min\_FOS}         & Global minimum FOS. \\
\texttt{FOS.critical\_member} & Element ID with lowest FOS. \\
\bottomrule
\end{tabularx}
\end{table}

\subsection{Utility Methods}

\begin{lstlisting}
analysis.resetAnalysis()
coords = analysis.getDeformedCoordinates(scale)
analysis.plotConvergence()
\end{lstlisting}

% -----------------------------------------------------------------------------
\section{\texttt{LoadCase}}
\label{sec:api:loadcase}
% -----------------------------------------------------------------------------

Manages nodal loads and load-factor scaling.

\begin{lstlisting}
lc = LoadCase(geometry);
lc.addNodalLoad(node_id, fx, fy, mz, 'description')
lc.setLoadFactor(1.0)
lc.clearLoads()

P  = lc.getScaledLoad()          % [Dx1] current load vector
Pf = lc.getFreeDOFLoads()        % [Fx1] free-DOF loads
lc.applyToState(state)           % write loads into an AnalysisState

factors = lc.generateLoadSteps(10, 'linear')
lc.displayLoads()
\end{lstlisting}

% -----------------------------------------------------------------------------
\section{\texttt{StructurePlotter}}
\label{sec:api:plotter}
% -----------------------------------------------------------------------------

Visualisation utilities.  Elements are drawn as filled rectangles matching
their physical thickness.

\begin{lstlisting}
plotter = StructurePlotter(geometry)
plotter = StructurePlotter(geometry, analysis)

% High-level plots
plotter.plotStructure(opts)              % undeformed model
plotter.plotDeformedShape(scale, opts)   % original + deformed
plotter.plotStressDistribution(type)     % 'axial','bending','shear','combined'
plotter.plotLoads(load_case, scale)      % force arrows

% Low-level drawing (used by demo scripts)
plotter.Plot_Beam(x, y, thickness, color)
plotter.Plot_Rectangle(x, y, thickness, color)
plotter.plotSupports()
\end{lstlisting}

% -----------------------------------------------------------------------------
\section{\texttt{optimizeThickness}}
\label{sec:api:optimize}
% -----------------------------------------------------------------------------

Standalone function implementing stress-ratio thickness optimisation.

\begin{lstlisting}
[geometry, results] = optimizeThickness(geometry, analysis, load_case, opts)
\end{lstlisting}

\begin{warnbox}[In-place mutation]
The \texttt{geometry} object is modified in place.  The returned
\texttt{geometry} is the same handle.  If you need to preserve the original,
create a copy before calling this function.
\end{warnbox}

\begin{table}[H]
\centering
\caption{\texttt{optimizeThickness} options.}
\begin{tabularx}{\textwidth}{llcX}
\toprule
\textbf{Option} & \textbf{Type} & \textbf{Default} & \textbf{Description} \\
\midrule
\texttt{alpha}            & double  & 0.4     & Stress-ratio exponent. \\
\texttt{max\_iterations}  & int     & 100     & Iteration cap. \\
\texttt{convergence\_tol} & double  & 0.01    & Max relative $\Delta t$. \\
\texttt{move\_limit}      & double  & 0.2     & Per-iteration change fraction. \\
\texttt{t\_min}           & double  & 0.001   & Minimum thickness (m). \\
\texttt{t\_max}           & double  & 0.020   & Maximum thickness (m). \\
\texttt{safety\_factor}   & double  & 1.5     & Divisor on $\sigma_{\text{ult}}$. \\
\texttt{use\_nonlinear}   & logical & \texttt{false} & Use large-deflection analysis. \\
\texttt{num\_load\_steps} & int     & 2       & Load steps for nonlinear. \\
\texttt{verbose}          & logical & \texttt{true}  & Print progress table. \\
\bottomrule
\end{tabularx}
\end{table}

% -----------------------------------------------------------------------------
\section{\texttt{subdivideTruss}}
\label{sec:api:subdivide}
% -----------------------------------------------------------------------------

\begin{lstlisting}
[new_nodes, new_conn, member_map] = subdivideTruss(nodes, conn, num_sub)
\end{lstlisting}

Replaces every member with a chain of \texttt{num\_sub} smaller elements.
Original joint nodes are preserved; intermediate nodes are inserted at
uniform spacing along each member.

% -----------------------------------------------------------------------------
\section{Importers}
\label{sec:api:importers}
% -----------------------------------------------------------------------------

\begin{lstlisting}
result = importCSVCurve(filename, opts)
result = importIGSCurve(filename, opts)
\end{lstlisting}

Both return a struct with fields \texttt{nodes}, \texttt{connectivity},
\texttt{num\_nodes}, \texttt{num\_elements}, \texttt{arc\_length},
\texttt{start\_tangent}, \texttt{end\_tangent}, \texttt{start\_angle}, and
\texttt{end\_angle}.

Common options: \texttt{num\_elements} (default 250), \texttt{scale}
(default~1), \texttt{flip} (default \texttt{false}), \texttt{plot} (default
\texttt{true}), \texttt{verbose} (default \texttt{true}).

% -----------------------------------------------------------------------------
\section{\texttt{exportProfile}}
\label{sec:api:export}
% -----------------------------------------------------------------------------

\begin{lstlisting}
profile = exportProfile(geometry, opts)
\end{lstlisting}

Rasterises the structure into a binary image, extracts boundary contours
using \texttt{bwboundaries}, fits smoothing splines, and exports each loop
as a SolidWorks-compatible text file.  See \Cref{ch:export} for a detailed
description of the algorithm.

\begin{table}[H]
\centering
\caption{\texttt{exportProfile} options.}
\begin{tabularx}{\textwidth}{llcX}
\toprule
\textbf{Option} & \textbf{Type} & \textbf{Default} & \textbf{Description} \\
\midrule
\texttt{num\_spline\_points}  & int    & 1000          & Points per loop. \\
\texttt{smoothing}            & double & 0.999999999   & \texttt{csaps} smoothing factor. \\
\texttt{scale\_to\_mm}        & double & 1000          & Coordinate multiplier for export. \\
\texttt{resolution}           & double & 50000         & Pixels per metre. \\
\texttt{min\_thickness\_frac} & double & 0             & Filter thin elements. \\
\texttt{export\_filename}     & char   & \texttt{[]}   & Base name for output files. \\
\texttt{plot}                 & logical& \texttt{true}  & Show diagnostic figure. \\
\texttt{verbose}              & logical& \texttt{true}  & Print summary. \\
\bottomrule
\end{tabularx}
\end{table}